\documentclass[a4paper, 12pt]{report}
\usepackage{fontspec}
\setmainfont{Times New Roman}
% \usepackage[utf8]{vietnam,inputenc}
% \usepackage[utf8]{inputenc}
% \usepackage[T5]{fontenc}


\DeclareTextSymbolDefault{\DH}{T1}

\usepackage{float}



\usepackage{nomencl,etoolbox,ragged2e,siunitx}
\usepackage{indentfirst, caption}
\usepackage{listings}
\usepackage{subcaption}
% \usepackage[final]{graphicx}
\usepackage{graphicx}
\usepackage[export]{adjustbox}
\usepackage[unicode,breaklinks]{hyperref}
\hypersetup{
    urlcolor=blue,
    linkcolor=black,
    citecolor=blue,
    colorlinks=true
} 
\usepackage[acronym]{glossaries}


% table vertical
\usepackage{lscape}

\renewcommand{\chaptername}{Chương}
\renewcommand{\contentsname}{Mục lục}
\renewcommand{\acronymname}{Danh mục viết tắt}
\makenomenclature
\usepackage{a4wide,amssymb,epsfig,latexsym,array,float,hhline,fancyhdr}
\usepackage[normalem]{ulem}
\usepackage[makeroom]{cancel}
\usepackage{amsmath}
\usepackage{amsthm}
\usepackage{multicol,longtable,amscd}
\usepackage{diagbox}%Make diagonal lines in tables
\usepackage[table,xcdraw,dvipsnames]{xcolor}
\usepackage{booktabs}
\usepackage{tikz}
\usepackage{pgf-pie}
\usepackage{rotating}
\usepackage{alltt}
\usepackage[framemethod=tikz]{mdframed}% For highlighting paragraph backgrounds
\usepackage{caption,subcaption}
\usepackage{tablefootnote}

\usepackage{lastpage}
\usepackage{enumitem}
\usepackage{color}
\usepackage{graphicx}							% Standard graphics package
\usepackage{array}
\usepackage{tabularx, caption}
\usepackage{multirow}
\usepackage{multicol}
\usepackage{rotating}
\usepackage{graphics}
\usepackage{geometry}
\usepackage{setspace}
\usepackage{epsfig}
\usepackage{tikz}
\usetikzlibrary{arrows,backgrounds}
% \usepackage[ruled,vlined]{algorithm2e}
\usepackage{subfiles}
\usepackage{algorithm}
\usepackage{algpseudocode}


%%%%% CROSSREF BIB
% \def\biblio{\bibliographystyle{amsalpha}\bibliography{bibgraf}\textbf{}} 

%\usepackage{pstcol} 								% PSTricks with the standard color package

\AtBeginDocument{\addtocontents{toc}{\protect\thispagestyle{empty}}} 

\AtBeginDocument{%
  \renewcommand\lstlistlistingname{Danh sách lệnh}
  \renewcommand\lstlistingname{Mã lệnh}
  
  \renewcommand\listtablename{Danh sách bảng}
  \renewcommand\tablename{Bảng}

  
\renewcommand{\listfigurename}{Danh sách hình vẽ}
\renewcommand{\figurename}{Hình}
}

\usepackage[normalem]{ulem}
\patchcmd{\chapter}{plain}{empty}{}{}
\patchcmd{\part}{plain}{empty}{}{}

\def\thesislayout{% A4: 210 × 297
	\geometry{
		a4paper,
		% total={160mm,240mm},  % fix over page
		left=30mm,
		right=20mm,
		top=25mm,
		bottom=20mm
	}
}
\thesislayout

\usepackage{fancyvrb}
\setlength{\headheight}{40pt}
\pagestyle{fancy}
\fancyhead{} % clear all header fields

\fancyhead[L]{
    \begin{tabular}{rl}
        \begin{picture}(10,0)(0,8)
            \includegraphics[width=8mm, height=8mm]{images/logoBK.jpg}
        \end{picture}&
        
        \begin{tabular}{l}
    	Đại học Quốc gia TP.HCM - Trường Đại Học Bách Khoa\\
    	Khoa Khoa học và Kỹ thuật Máy tính
        \end{tabular} 	
    \end{tabular}
}
\fancyfoot{} % clear all footer fields
\fancyfoot[L]{\footnotesize Đồ án tốt nghiệp}
\fancyfoot[R]{\footnotesize Trang {\thepage}/\pageref{lastPage}}
\renewcommand{\headrulewidth}{0.3pt}
\renewcommand{\footrulewidth}{0.3pt}


%%%
\setcounter{secnumdepth}{4}
\makeatletter
\newcounter {subsubsubsection}[subsubsection]
\renewcommand\thesubsubsubsection{\thesubsubsection .\@alph\c@subsubsubsection}
\newcommand\subsubsubsection{\@startsection{subsubsubsection}{4}{\z@}%
                                     {-3.25ex\@plus -1ex \@minus -.2ex}%
                                     {1.5ex \@plus .2ex}%
                                     {\normalfont\normalsize\bfseries}}
\newcommand*\l@subsubsubsection{\@dottedtocline{3}{10.0em}{4.1em}}
\newcommand*{\subsubsubsectionmark}[1]{}
\def\toclevel@subsubsubsection{5}
\makeatother

% \everymath{\color{blue}}%make in-line maths symbols blue to read/check easily

\sloppy
\captionsetup[figure]{labelfont={small,bf},textfont={small,it},belowskip=-1pt,aboveskip=-9pt}
%space remove between caption, figure, and text
\captionsetup[table]{labelfont={small,bf},textfont={small,it},belowskip=-1pt,aboveskip=7pt}



\setlength{\floatsep}{5pt plus 2pt minus 2pt}
\setlength{\textfloatsep}{5pt plus 2pt minus 2pt}
\setlength{\intextsep}{10pt plus 2pt minus 2pt}

\thesislayout
%code
\usepackage{color} % tô màu cho code
\definecolor{dkgreen}{rgb}{0,0.6,0}
\definecolor{gray}{rgb}{0.5,0.5,0.5}
\definecolor{mauve}{rgb}{0.58,0,0.82}

\definecolor{codegreen}{rgb}{0,0.6,0}
\definecolor{codegray}{rgb}{0.5,0.5,0.5}
\definecolor{codepurple}{rgb}{0.58,0,0.82}
\definecolor{backcolour}{rgb}{0.95,0.95,0.92}

\lstset{
language=XML,
    aboveskip=3mm,
    belowskip=3mm,
    columns=flexible,
    backgroundcolor=\color{backcolour},   
    commentstyle=\color{codegreen},
    keywordstyle=\color{magenta},
    numberstyle=\tiny\color{codegray},
    stringstyle=\color{codepurple},
    basicstyle=\ttfamily\footnotesize,
    breakatwhitespace=false,         
    breaklines=true,                 
    captionpos=b,                    
    keepspaces=true,                 
    numbers=left,                    
    numbersep=5pt,                  
    showspaces=false,                
    showstringspaces=false,
    showtabs=false,                  
    tabsize=2,
    otherkeywords = {!,!=,~,\$,*,\&,\%/\%,\%*\%,\%\%, \%in\%,<-,<<-,/},
    alsoletter={.,_}
}



% Glossary
\makeglossaries
% \newacronym{gcd}{GCD}{Greatest Common Divisor}
% \newacronym{lcm}{LCM}{Least Common Multiple}
% \newacronym{bkel}{BKEL}{Bách Khoa e-Learning}
% \newacronym{ags}{AGS}{Auto Grading System}
\newacronym{pgsts}{PGS.TS}{Phó giáo sư tiến sĩ}
\newacronym{ths}{ThS}{Thạc sĩ}
\newacronym{dhbk}{ĐHBK - ĐHQG-HCM}{Trường Đại học Bách khoa - Đại học
Quốc gia Thành phố Hồ Chí Minh}
\newacronym{dclv}{ĐCLV}{Đề cương luận văn}
\newacronym{dacn}{ĐACN}{Đồ án chuyên ngành}
\newacronym{dss}{DSS}{Hệ thống hỗ trợ quyết định (Decision Support System)}
\newacronym{mcdm}{MCDM}{Phân tích Quyết định Đa tiêu chí (Multiple-criteria decision-making)}
\newacronym{mcda}{MCDA}{Phân tích Quyết định Đa tiêu chí (Multi-Criteria Decision Analysis)}
\newacronym{mbti}{MBTI}{Myers-Briggs Type Indicator}
% \newacronym{db}{CSDL}{Cơ sở dữ liệu}
% \newacronym{cntt}{CNTT}{Công nghệ thông tin}
% \newacronym{fp}{KTLT}{Kỹ thuật lập trình}
% \newacronym{dsa}{CTDL\&GT}{Cấu trúc dữ liệu và Giải thuật}
% \newacronym{ui}{UI}{User Interface}
% \newacronym{html}{HTML}{HyperText Markup Language}
% \newacronym{css}{CSS}{Cascading Style Sheets}
% \newacronym{api}{API}{Application Programming Interface}
% \newacronym{http}{HTTP}{Hypertext Transfer Protocol}
% \newacronym{rest}{REST}{Representational state transfer}
% \newacronym{bi}{BI}{Business Intelligence}
% \newacronym{ai}{AI}{Artificial Intelligence}
% \newacronym{ml}{ML}{Machine Learning}
% \newacronym{hdfs}{HDFS}{Hadoop Distributed File System}
% \newacronym{acid}{ACID}{Atomicity, Consistency, Isolation, and Durability}
% \newacronym{dbms}{DBMS}{Database Management System}
% \newacronym{rdbms}{RDBMS}{Relational Database Management System}
% \newacronym{rpc}{RPC}{Remote Procedure Call}
% \newacronym{bson}{BSON}{Binary Javascript Object Notation}
% \newacronym{json}{JSON}{Javascript Object Notation}
% \newacronym{csdl}{CSDL}{Cơ sở dữ liệu}
% \newacronym{olap}{OLAP}{Online Analytical Processing}
% \newacronym{oltp}{OLTP}{Online Transaction Processing}
% \newacronym{iot}{IoT}{Internet of Things}
% \newacronym{aws}{AWS}{Amazon Web Services}
% \newacronym{sql}{SQL}{Structured Query Language}
% \newacronym{sla}{SLA}{Service Level Agreement}
% \newacronym{csv}{CSV}{Comma Separated Values}
% \newacronym{etl}{ETL}{Extract - Transform - Load}
% \newacronym{crm}{CRM}{Customer Relationship Management}
% \newacronym{nosql}{NoSQL}{Not Only SQL}
% \newacronym{saas}{SaaS}{Software as a Service}
% \newacronym{kdc}{KDC}{Key distribution center}
% \newacronym{tgt}{TGT}{Ticket Granting Tickets}
% \newacronym{jwt}{JWT}{JSON Web Token}
% \newacronym{rsa}{RSA}{Ron Rivest, Adi Shamir and Leonard Adleman}
% \newacronym{ecdsa}{ECDSA}{Elliptic Curve Digital Signature Algorithm}
% \newacronym{hmac}{HMAC}{Hashed Message Authentication Code}
% \newacronym{sha256}{SHA256}{Secure Hash Algorithm 256-bit}
% \newacronym{io}{I/O}{Input/Output}
% \newacronym{cron}{Cron}{Command Run On}
% \newacronym{rdd}{RDD}{Resilient Distributed Dataset}
% \newacronym{dom}{DOM}{Document Object Model}
% \newacronym{mvc}{MVC}{Model - View - Controller}
% \newacronym{xlsx}{XLSX}{Excel Microsoft Office Open XML Format Spreadsheet file}
% \newacronym{ip}{IP}{Internet Protocol address}
% \newacronym{ram}{RAM}{Random Access Memory}
% \newacronym{yarn}{YARN}{Yet Another Resource Negotiator}
% \newacronym{ssh}{SSH}{Secure Socket Shell}
% \newacronym{jdbc}{JDBC}{Java database connectivity}
% \newacronym{erd}{ERD}{Entity Relationship Diagram}
% \newacronym{des}{DES}{Data Encryption Standard}
% \newacronym{rc4}{RC4}{Rivest Cipher 4}
% \newacronym{aes}{AES}{Advanced Encryption Standard}
% \newacronym{rfc}{RFC}{Request for Comments}
% \newacronym{cpu}{CPU}{Central Processing Unit}
% \newacronym{dsl}{DSL}{Domain-specific Language}
\setstretch{1.25}





\begin{document}

\begin{titlepage}
    \begin{tikzpicture}[remember picture,overlay,inner sep=0,outer sep=0]
         \draw[black,line width=0.5pt] ([xshift=-1.5cm,yshift=-2cm]current page.north east) coordinate (A)--([xshift= 2.5cm,yshift=-2cm]current page.north west) coordinate(B)--([xshift=2.5cm,yshift=2cm]current page.south west) coordinate (C)--([xshift=-1.5cm,yshift=2cm]current page.south east) coordinate(D)--cycle;
         
         \draw[black,line width=4pt] ([xshift=-1.6cm,yshift=-2.1cm]current page.north east) coordinate (A)--([xshift= 2.6cm,yshift=-2.1cm]current page.north west) coordinate(B)--([xshift= 2.61cm,yshift=2.1cm]current page.south west) coordinate (C)--([xshift=-1.6cm,yshift=2.1cm]current page.south east) coordinate(D)--cycle;
    \end{tikzpicture}
    
    \begin{center}
        \fontsize{15pt}{15pt}\selectfont
        \textbf{ĐẠI HỌC QUỐC GIA TP.HCM}\\
        \textbf{TRƯỜNG ĐẠI HỌC BÁCH KHOA}\\
        \textbf{KHOA KHOA HỌC VÀ KỸ THUẬT MÁY TÍNH}
    \end{center}
    
    \begin{figure}[h!]
        \centering
        \includegraphics[width=4cm]{images/logoBK.jpg}
    \end{figure}%
    
    \begin{tikzpicture}
        \node at (.45\textwidth, 2) [name=DACN]{\fontsize{15}{15}\selectfont\textbf{BÁO CÁO}};
        \node at (.45\textwidth, 1) [name=CO2001]{\fontsize{15}{15}\selectfont\textbf{ĐỒ ÁN TỐT NGHIỆP}};
        \draw[thick,white] (-.5,.5) -- (.98\textwidth,.5);
        \node at (.5,0){\fontsize{14}{14}\selectfont\textbf{}};
        \node at (.475\textwidth, -1){\fontsize{18}{20}\selectfont\textbf{XÂY DỰNG HỆ THỐNG}};
        \node at (.475\textwidth, -2){\fontsize{18}{20}\selectfont\textbf{HỖ TRỢ QUYẾT ĐỊNH}};
        \node at (.475\textwidth, -3){\fontsize{18}{20}\selectfont\textbf{CHỌN NGÀNH HỌC ĐẠI HỌC}};
        \node at (.475\textwidth, -4){\fontsize{15}{16}\selectfont\textbf{NGÀNH: KHOA HỌC MÁY TÍNH}};
        \draw[thick,white] (-.5,-2.7) -- (.98\textwidth,-2.7);
    \end{tikzpicture}\par%5
    
    \vspace{1.5cm}
    \begin{table}[h!]
        \raggedleft
        \bfseries
        \renewcommand{\arraystretch}{1.2}
        \fontsize{15pt}{15pt}\selectfont
        \begin{tabular}{llc}
            HỘI ĐỒNG:   & Hội đồng X Khoa học Máy tính &\\
            GVHD 1:     & ThS. Bùi Tiến Đức \\
            GVHD 2:     & PGS.TS Trần Minh Quang \\
            % GVPB:       & ... & \\
            
            & \quad\quad---o0o--- &  \\
            SVTH 1:     & Nguyễn Quang Minh      & 2111753 \\
            SVTH 2:     & Nguyễn Doãn Hoàng      & 2111238 \\
            SVTH 3:     & Nguyễn Trần Ngọc Liên  & 2011521 \\
        \end{tabular}
    \end{table}
    
    \begin{center}
        \vspace{1cm}\fontsize{14pt}{9pt}\selectfont
        
        TP. HỒ CHÍ MINH, 05/2024
    \end{center}
\end{titlepage}


%%%%% MARKUP BEGIN CHUONG 0 %%%%%%%%%%%%%
\pagestyle{empty}
% TODO: đánh số i ii
\pagenumbering{roman}



\chapter*{Lời cam đoan}
Chúng tôi xin cam đoan Đồ án tốt nghiệp này là công trình nghiên cứu khoa học nghiêm túc và trung thực của nhóm dưới sự hướng dẫn của \acrshort{pgsts} Trần Minh Quang và \acrshort{ths}. Bùi Tiến Đức. Tất cả nội dung được trình bày trong Đồ án tốt nghiệp này, ngoại trừ những phần đã được chú thích, trích nguồn rõ ràng trong phần tài liệu tham khảo, đều là do chính bản thân nhóm thực hiện.

\indent Chúng tôi xin chịu trách nhiệm hoàn toàn nếu có bất kỳ sự gian lận nào về nội dung Đồ án tốt nghiệp của nhóm.

\begin{flushright}
    \textbf{Nhóm thực hiện}\\
    Nguyễn Quang Minh \\
    Nguyễn Doãn Hoàng \\
    Nguyễn Trần Ngọc Liên \\
\end{flushright}

\chapter*{Lời cảm ơn}
Đồ án tốt nghiệp được hoàn thành dưới sự hướng dẫn  của PGS.TS Trần Minh Quang và \acrshort{ths}. Bùi Tiến Đức, Khoa Khoa học và Kỹ thuật Máy tính, Trường Đại học Bách khoa - Đại học Quốc gia Thành phố Hồ Chí Minh. Nhóm thực hiện xin chân thành cảm ơn PGS.TS Trần Minh Quang và \acrshort{ths}. Bùi Tiến Đức, đã giúp đỡ chúng em kiến thức chuyên môn, thảo luận, đưa ra gợi ý và tạo điều kiện thuận lợi cho chúng em hoàn thành đồ án.

% Trong suốt quá trình nghiên cứu đề tài, nhóm chúng em cũng đã nhận được nhiều sự giúp đỡ về hạ tầng hệ thống từ thầy Nguyễn Mạnh Thìn, anh Nguyễn Đức Hiệp, anh Trần Hữu Lý, anh Nguyễn Văn Quyền Lực và giải đáp thắc mắc từ nhóm luận văn đi trước. Nhóm em xin cảm ơn sự giúp đỡ nhiệt tình từ các thầy, các anh để chúng em hiện thực thành công đề tài này.


Chúng em cũng xin gửi lời cảm ơn đặc biệt đến quý thầy cô phản biện, những người đã đọc và đóng góp ý kiến để chúng em hoàn thiện Đồ án tốt nghiệp của mình.


Nhóm thực hiện cũng xin bày tỏ lòng biết ơn sâu sắc đến quý thầy cô Khoa Khoa học và Kỹ thuật Máy tính, Trường Đại học Bách khoa - Đại học Quốc gia Thành phố Hồ Chí Minh, là những người đã truyền thụ kiến thức chuyên môn, đã tạo điều kiện cho chúng em học tập và phát triển trong suốt quãng thời gian vừa qua.

Kính chúc quý thầy cô sức khoẻ, thành công và tiếp tục đào tạo những thế hệ sinh viên mới trong tương lai.


Chúng em xin chân thành cảm ơn.



\begin{flushright}
    \textbf{Nhóm thực hiện}\\
    Nguyễn Quang Minh \\
    Nguyễn Doãn Hoàng \\
    Nguyễn Trần Ngọc Liên \\
\end{flushright}

\chapter*{Tóm tắt đề tài}
Nghiên cứu và phát triển Hệ thống hỗ trợ quyết định (Decision Support System - \acrshort{dss}) bao gồm hệ thống website lấy ý tưởng từ Ikigai - Lý do để sống của người Nhật nhằm hỗ trợ học sinh lựa chọn chuyên ngành học thuật, ngành nghề tương lai. Mục tiêu chính của hệ thống là giúp học sinh khám phá ra niềm đam mê, sở trường của mình và cách chúng phù hợp với nhu cầu của xã hội, từ đó thúc đẩy các bạn học sinh đưa ra những quyết định hợp lý về con đường học vấn và nghề nghiệp.

Phương pháp được thực hiện trên hệ thống bao gồm việc mô hình hóa bài toán sử dụng sơ đồ Ikigai Venn với 4 câu hỏi: ``Bạn yêu thích điều gì?", ``Bạn giỏi gì?", ``Bạn có thể được trả tiền để làm gì?" và ``Thế giới cần gì?". Việc giải quyết được thực hiện bằng cách trả lời cho từng câu hỏi, sau đó tổng hợp và tính toán kết quả bằng các phương pháp giải quyết vấn đề phù hợp dựa trên Phân tích Quyết định Đa tiêu chí (Multi-Criteria Decision Making - \acrshort{mcdm}). Phương pháp \acrfull{mbti} được sử dụng để tìm hiểu tính cách của người dùng trả lời cho câu hỏi ``Bạn yêu thích điều gì?", trong khi Nhóm ngành nghề (Career Clustering) giúp người dùng xác đinh nhóm ngành dựa trên thế mạnh trả lời cho câu hỏi ``Bạn giỏi gì?". Dữ liệu từ nhiều tổ chức khác nhau như tổng cục thống kê, phòng thống kê của tổ chức lao động thế giới ILOSTAT và nhiều nguồn uy tín khác được tổng hợp để trả lời các câu hỏi liên quan đến ``Bạn có thể được trả tiền để làm gì?" và ``Thế giới cần gì?". Các giải thuật về Weight-sum và Vikor trong bài toán phân tích quyết định đa tiêu chí được sử dung để tính toán giải pháp tốt nhất.

Dưới sự gợi ý của hệ thống hỗ trợ quyết định (\acrshort{dss}), học sinh sẽ có được hiểu biết sâu sắc hơn về đam mê, sở trường và các con đường nghề nghiệp tiềm năng, từ đó dẫn đến sự tự tin và rõ ràng hơn trong quá trình ra quyết định. Khảo sát sơ bộ trên 100 sinh viên mang lại kết quả tích cực, cho thấy sự hài lòng của học sinh về hệ thống, cũng như độ chính xác ở mức tốt. Những phát hiện này cho thấy hiệu quả của phương pháp được triển khai và mở đường trong tương lai cho việc đánh giá và tinh chỉnh thêm để liên tục nâng cao khả năng của hệ thống trong việc hỗ trợ học sinh trên hành trình học tập và nghề nghiệp.

% Mục lục
\tableofcontents

% Danh sách bảng
\listoftables%[title={Danh sách bảng}]

% Danh sách hình vẽ
\listoffigures%[til.101 \msg_fatal:nn {fontspec} {cannot-use-pdftex}tle={Danh sách hình vẽ}]

% Danh sách code
% \lstlistoflistings%[title={Danh sách lệnh}]

% Danh mục viết tắt
\printglossary[type=\acronymtype,title={Danh mục viết tắt}]
%%%%% MARKUP END CHUONG 0 %%%%%%%%%%%%%

\newpage
\pagestyle{fancy}
\pagenumbering{arabic}

%%%%% MARKUP BEGIN CHUONG 1 %%%%%%%%%%%%%
\chapter{Tổng quan về đề tài}\label{chap:overview}
% Cần phải set riêng cho từng chương vì đã set page style cho chapter là empty ở \patchcmd{\chapter}{plain}{empty}{}{}
\thispagestyle{fancy} % show footer for chaper page

\textit{Nội dung chương 1 trình bày sơ lược về đề tài, bối cảnh hiện nay và làm rõ lý do đề tài được chọn. Từ đó, xác định mục tiêu cần đạt được khi thực hiện và phạm vi của đề tài.}

\import{chapters/1. Tổng quan/}{main.tex}
%%%%% MARKUP END CHUONG 1 %%%%%%%%%%%%%

%%%%% MARKUP BEGIN CHUONG 2 %%%%%%%%%%%%%
\chapter{Nền tảng lý thuyết (Background) và Các nghiên cứu liên quan (Related work)}\label{chap:chap2}
\thispagestyle{fancy} % show footer for chaper page

\textit{Chương 2 đề cập tới các khái niệm và nền tảng lý thuyết được sử dụng để xây dựng hệ thống và phân tích các nghiên cứu liên quan.}

\import{chapters/2. Nền tảng lý thuyết (Background) và Các nghiên cứu liên quan (Related work)/}{main.tex}
%%%%% MARKUP END CHUONG 2 %%%%%%%%%%%%%

%%%%% MARKUP BEGIN CHUONG 3 %%%%%%%%%%%%%
\chapter{Phương pháp thực hiện}\label{chap:chap3}
\thispagestyle{fancy} % show footer for chaper page

\textit{Chương 3 .}

\import{chapters/3. Phương pháp thực hiện/}{main.tex}
%%%%% MARKUP END CHUONG 3 %%%%%%%%%%%%%

%%%%% MARKUP BEGIN CHUONG 4 %%%%%%%%%%%%%
\chapter{Thiết kế hệ thống}\label{chap:chap4}
\thispagestyle{fancy} % show footer for chaper page

\textit{Chương 4 .}

\import{chapters/4. Thiết kế hệ thống/}{main.tex}
%%%%% MARKUP END CHUONG 3 %%%%%%%%%%%%%

%%%%% MARKUP BEGIN CHUONG 5 %%%%%%%%%%%%%
\chapter{Hiện thực hệ thống}\label{chap:chap5}
\thispagestyle{fancy} % show footer for chaper page

\textit{Chương 5 .}

\import{chapters/5. Hiện thực hệ thống/}{main.tex}
%%%%% MARKUP END CHUONG 5 %%%%%%%%%%%%%

%%%%% MARKUP BEGIN CHUONG 6 %%%%%%%%%%%%%
\chapter{Đánh giá kết quả}\label{chap:chap6}
\thispagestyle{fancy} % show footer for chaper page

\textit{Chương 6 thể hiện quá trình tiến hành đánh giá kết quả dựa trên việc khảo sát trực tiếp ngẫu nhiên 100 bạn sinh viên tại trường đại học Bách Khoa và thu được những kết quả bước đầu hết sức khả quan.}

\import{chapters/6. Đánh giá kết quả/}{main.tex}
%%%%% MARKUP END CHUONG 6 %%%%%%%%%%%%%

%%%%% MARKUP BEGIN CHUONG 7 %%%%%%%%%%%%%
\chapter{Tổng kết} \label{chap:summary}
\thispagestyle{fancy} % show footer for chaper page

\textit{Nội dung chương \ref{chap:summary} trình bày đánh giá về quá trình hiện thực đồ án tốt nghiệp đối với các mục tiêu đề ra và đánh giá ưu, nhược điểm của hệ thống đã xây dựng. Từ đó thảo luận các hướng phát triển trong tương lai tiếp tục để cải thiện.}

\import{chapters/7. Tổng kết/}{main.tex}
%%%%% MARKUP END CHUONG 7 %%%%%%%%%%%%%

%%%%% MARKUP BEGIN PHU LUC %%%%%%%%%%%%%
\newpage
\pagestyle{empty}

\import{chapters/Tài liệu tham khảo/}{main.tex}

% \newpage
% \pagestyle{fancy}

\end{document}